% ─────────────────────────────────────────────────────────────────────────────
\section{Methodology}
\label{sec:methodology}
% ─────────────────────────────────────────────────────────────────────────────

\subsection{Problem Formulation}
\label{subsec:formulation}

Given a flight represented as a feature vector
$\mathbf{f} \in \mathbb{R}^d$ (from any representation in
Section~\ref{sec:features}), the system must produce one of 20 diagnostic labels:

\begin{equation}
    \hat{y} \in \{0, 1, 2, \ldots, 19\}
\end{equation}

where label 0 denotes a \textit{healthy} flight and labels 1--19 denote the
19 fault categories of Table~\ref{tab:class_dist}. This 20-class problem is
decomposed into two sequential sub-tasks:

\begin{enumerate}
    \item \textbf{Anomaly Detection (AD):} Binary classification,
          $y_{\text{AD}} \in \{0 = \text{healthy},\ 1 = \text{anomaly}\}$.
    \item \textbf{Fault Classification (FC):} 19-class classification,
          $y_{\text{FC}} \in \{1, 2, \ldots, 19\}$, performed \textit{only} on
          flights flagged as anomalous by the AD stage.
\end{enumerate}

\subsection{Two-Stage Diagnosis Pipeline}
\label{subsec:pipeline}

The two-stage cascade architecture is motivated by the observation that AD and
FC require fundamentally different inductive biases: AD benefits from global
whole-flight context, while FC requires local fault-specific pattern resolution
\cite{lmsd2025}. Jointly optimizing both objectives in a single model leads to
a conflict where the global context needed for AD suppresses the local fault
signatures needed for FC.

The pipeline routes each incoming flight through Stage~1 (AD). Only flights
classified as anomalous proceed to Stage~2 (FC). Healthy flights receive label 0
and bypass the fault classifier entirely.

Figure~\ref{fig:pipeline} illustrates the architecture.

\begin{figure}[H]
    \centering
    \begin{tikzpicture}[
        every node/.style={align=center, font=\small},
        box/.style={draw, rounded corners=3pt, minimum width=2.2cm,
                    minimum height=0.8cm, inner sep=5pt},
        arr/.style={-{Stealth[length=5pt]}, thick},
        node distance=0.7cm and 1.3cm
    ]
        \node[box, fill=gray!15]   (raw)  {Flight\\(features $\mathbf{f}$)};
        \node[box, fill=blue!10, right=of raw]  (ad)   {\textbf{Stage 1}\\AD model};
        \node[box, fill=green!15, above right=0.4cm and 1.2cm of ad] (h)
              {Label 0\\HEALTHY};
        \node[box, fill=orange!15, below right=0.4cm and 1.2cm of ad] (fc)
              {\textbf{Stage 2}\\FC model};
        \node[box, fill=red!12, right=of fc] (fault) {Labels 1--19\\fault type};
        \draw[arr] (raw) -- (ad);
        \draw[arr] (ad)  -- node[above,font=\scriptsize,sloped]{healthy} (h);
        \draw[arr] (ad)  -- node[below,font=\scriptsize,sloped]{anomaly} (fc);
        \draw[arr] (fc)  -- (fault);
    \end{tikzpicture}
    \caption{Two-stage diagnosis pipeline. Stage~1 (AD) routes flights to the
             healthy label or to Stage~2 (FC) for fault type identification.}
    \label{fig:pipeline}
\end{figure}

\subsection{Classical Models}
\label{subsec:classical}

Five classical scikit-learn models are evaluated on both the AD and FC tasks:

\paragraph{Logistic Regression (LR).}
A linear probabilistic classifier trained by maximizing the regularized
log-likelihood. The decision boundary is a hyperplane in feature space:
$\hat{p}(y=1 \mid \mathbf{f}) = \sigma(\mathbf{w}^\top \mathbf{f} + b)$,
where $\sigma$ is the sigmoid function. LR is interpretable, fast to train,
and robust to high-dimensional sparse features. The L2 regularization strength
$C$ is tuned per fold.

\paragraph{Linear Support Vector Machine (LinearSVM).}
Finds the maximum-margin hyperplane separating the two classes. Under class
imbalance, the \texttt{class\_weight=`balanced'} option rescales the hinge loss
by inverse class frequency. LinearSVM is equivalent to LR in the linear limit but
optimizes the margin directly rather than the log-likelihood.

\paragraph{XGBoost (XGB).}
An ensemble of gradient-boosted decision trees \cite{xgboost2016}. XGBoost
captures nonlinear feature interactions and implicitly handles feature
importance through its information-gain-based splitting criterion. The
\texttt{scale\_pos\_weight} parameter balances positive and negative class
weights for AD; \texttt{class\_weight=`balanced'} is used for FC.

\paragraph{Random Forest (RF).}
A bagged ensemble of decision trees with feature subsampling at each split.
RF provides implicit feature selection through out-of-bag importance scores
and is robust to irrelevant features---important given that Repr.~B contains
184 dimensions of which many are low-information for specific fault types.

\paragraph{K-Nearest Neighbors (KNN).}
Distance-based classifier that assigns the majority label among the $k$ nearest
neighbors in feature space. KNN serves as a sanity-check baseline; its
performance is sensitive to the feature scale and to the curse of dimensionality
in high-dimensional representations.

\subsection{Deep Learning Benchmarks}
\label{subsec:dl}

Three deep learning architectures from the NGAFID benchmark paper
\cite{convtok2023, lmsd2025} serve as upper-bound comparisons. Unlike
classical models, these operate directly on the raw 2{,}048-step time series,
preserving all temporal resolution.

\paragraph{ConvTokMHSA.}
The ConvTok framework \cite{convtok2023} first tokenizes the multivariate time
series into non-overlapping patches using 1D CNN encoders that capture local
shape descriptors. The tokens are then processed by a Multi-Head Self-Attention
(MHSA) layer with full global receptive field---every token attends to every
other token across the entire flight. This global aggregation is optimal for AD,
where the overall health state is determined by the coherence of the whole-flight
profile (F1\,=\,0.799). However, it suffers on FC because cross-stage
operational noise dilutes weak local fault signatures (F1\,=\,0.208).

\paragraph{MMK Net.}
The Multi-Micro Kernel Network \cite{lmsd2025} employs parallel 1D convolutions
with micro-scale kernels ($k = 1, 3, 5$) that explicitly restrict the receptive
field to local temporal windows. This constraint prevents operational noise from
dominating and isolates micro-pattern fault signatures. MMK~Net deliberately
omits pooling layers to preserve temporal resolution, and uses LayerNorm instead
of BatchNorm to stabilize training under severe class imbalance. It achieves the
best FC performance (F1\,=\,0.623).

\paragraph{LMSD.}
The Long-Micro Scale Diagnostician \cite{lmsd2025} is the state-of-the-art
end-to-end architecture that explicitly operationalizes the two-stage decomposition
principle. It uses ConvTokMHSA (global receptive field) for Stage~1 AD and
MMK~Net (restricted receptive field) for Stage~2 FC, connected by a
\textit{hard-threshold routing mechanism} that enforces dimensional isolation:
healthy samples trigger a label-0 output with the fault dimensions masked to
$-\infty$, completely bypassing the FC head. Training is decoupled: the AD head
trains on all 11{,}446 flights, while the FC head trains exclusively on the
5{,}601 anomalous flights, preventing the large healthy class from suppressing
gradient signals for rare fault classes.

\subsection{Oracle Hierarchical Classification}
\label{subsec:oracle}

An additional experiment assesses the theoretical ceiling achievable through
physics-based fault grouping. Ward hierarchical clustering is applied to the
class-conditional feature centroids, grouping the 19 fault classes into $K$
physically coherent super-groups. A separate XGBoost classifier is then trained
for each group on the anomalous flights belonging to that group.

This \textit{Oracle} configuration uses ground-truth group labels at test time
(i.e., it knows which group a fault belongs to), providing an upper bound for
the hierarchical approach. The K\,=\,4 solution partitions the 19 classes as:

\begin{itemize}
    \item \textbf{Group 0} (156 flights): Classes 10, 12 --- cylinder structural failures
    \item \textbf{Group 1} (2{,}476 flights): Classes 6, 8, 14, 17 --- engine operational issues
    \item \textbf{Group 2} (2{,}875 flights): Classes 0--5, 7, 9, 11, 15, 16, 18 --- baffle/intake/engine maintenance
    \item \textbf{Group 3} (94 flights): Class 13 --- overspeed events
\end{itemize}

\subsection{Evaluation Metrics}
\label{subsec:metrics}

\subsubsection{F1-macro Score}

The primary metric is the \textbf{macro-averaged F1 score}:
\begin{equation}
    F1_{\text{macro}} = \frac{1}{K}\sum_{k=1}^{K} \frac{2 P_k R_k}{P_k + R_k}
\end{equation}
where $P_k$ and $R_k$ are precision and recall for class $k$, and $K$ is the
number of classes. Macro averaging treats all classes equally regardless of
support, making it sensitive to performance on minority fault classes---which
is precisely what matters for fault detection in long-tailed distributions.

\subsubsection{Area Under the ROC Curve (AUC)}

For AD, the AUC summarizes discrimination ability across all decision thresholds
and is reported alongside F1. AUC is threshold-independent and provides a
complementary view to F1 under varying operating points.

\subsubsection{False Negative Rate (FNR)}

The False Negative Rate measures the fraction of anomalous flights incorrectly
classified as healthy:
\begin{equation}
    \text{FNR} = \frac{FN}{FN + TP}
\end{equation}
In an aviation safety context, a false negative---a missed fault---is far more
consequential than a false positive (unnecessary inspection). FNR is therefore
a primary safety metric.

\subsubsection{MCWPM (Mission-Critical Weighted Performance Metric)}

MCWPM is a safety-weighted F-score that penalizes false negatives more heavily
than false positives:
\begin{equation}
    \text{MCWPM} = \frac{\alpha \cdot R + P}{\alpha + 1}
    \label{eq:mcwpm}
\end{equation}
where $P$ is precision, $R$ is recall, and $\alpha = 2.5$. Under this
parameterization, a false negative contributes $\alpha^2 = 6.25$ times more
to the loss than a false positive. MCWPM is reported alongside F1 and FNR for
all AD experiments to provide a safety-contextualized performance assessment.
